\documentclass[openany]{book}
\usepackage[utf8]{inputenc}
\usepackage{hyperref}
\usepackage{listings}
\usepackage{xcolor}
\usepackage{enumitem}
\usepackage[margin=1in]{geometry}

\lstset{
    basicstyle=\ttfamily\small,
    breaklines=true,
    frame=single,
    backgroundcolor=\color{gray!5},
}

\title{Consolidated Oneliner Library}
\author{GULECD-UFPR}
\date{\today}

\begin{document}
\maketitle
\tableofcontents


\chapter{Bash}

\section{Create Fallocate}
\section*{Create a large file of a specific size quickly}
\textbf{Author:} Marcos de Carvalho \hfill \textbf{Date:} 2026-04-21

\noindent Quickly creates a file of a specific size (e.g., 1 GB) using fallocate, which is much faster than dd because it allocates blocks without writing actual data.

\section*{Language}
bash

\section*{Category}
filesystem

\section*{Command}
\begin{lstlisting}
fallocate -l 1G test_file
\end{lstlisting}

\section*{Explanation}
The 'fallocate' command communicates directly with the filesystem to preallocate disk space for a file. Unlike 'dd', which must write every byte to disk, 'fallocate' simply marks the space as allocated in the filesystem's metadata. This is extremely efficient for creating large test files, swap files, or preallocating space for databases. The length (-l) can be specified in bytes or using suffixes like K, M, G, T, P, E.

\section*{Tags}
fallocate, filesystem, storage, disk-usage, allocation, testing

\section*{Dependencies}
util-linux


\section*{Arguments}
\begin{enumerate}
  \item \textbf{SIZE} (Optional): The size of the file to create (e.g., 10M, 1G, 500M). \\ Default: 1G
  \item \textbf{FILENAME} (Optional): The name of the file to be created. \\ Default: test\_file
\end{enumerate}

\section*{Examples}
\begin{enumerate}
  \item \texttt{fallocate -l 100M dummy\_file} - Create a 100 MB file named dummy\_file
  \item \texttt{fallocate -l 1G large\_test\_file} - Create a 1 GB file for testing
  \item \texttt{fallocate -l 10G /swapfile} - Preallocate 10 GB for a swap file (requires subsequent mkswap and swapon)
\end{enumerate}

\section*{Output}
\begin{lstlisting}
(No output on success)
\end{lstlisting}

\section*{Notes}
\begin{itemize}
  \item Not all filesystems support fallocate. It is well-supported on ext4, xfs, btrfs, and ocfs2.
  \item If the filesystem doesn't support fallocate, you can use 'truncate -s SIZE FILENAME' as a fallback, though it creates a sparse file.
  \item For older systems or filesystems, 'dd if=/dev/zero of=FILENAME bs=1M count=SIZE\_IN\_MB' is the most compatible but much slower method.
\end{itemize}

\section*{Warnings}
\begin{itemize}
  \item This command will overwrite the specified file if it already exists.
  \item On some filesystems, the allocated space may contain 'stale' data from previously deleted files, though modern Linux filesystems (ext4, xfs) zero-fill or mark it as uninitialized for security.
  \item Ensure you have enough free disk space before running this command.
\end{itemize}

\section*{See Also}
\begin{itemize}
  \item find-large-files-recursive
  \item disk-space-usage-per-directory
\end{itemize}

\section*{Status}
reviewed

\section*{Safety}
caution

\section*{Shell}
bash

\section*{Platforms}
linux

\section*{Created At}
2026-04-21

\section*{Updated At}
2026-04-21

\section{Disk Space Sort Largest Directories}
\section*{Sort directories by disk usage and display the largest ones}
\textbf{Author:} Marcos de Carvalho \hfill \textbf{Date:} 2026-04-21

\noindent Displays the top 10 directories with the largest disk usage in the current directory, sorted in descending order by size in human-readable format.

\section*{Language}
bash

\section*{Category}
disk-usage

\section*{Command}
\begin{lstlisting}
du -sh */ | sort -hr | head -10
\end{lstlisting}

\section*{Explanation}
The command uses 'du -sh' to calculate the total disk usage for each directory in human-readable format. The '*/' glob pattern matches only directories (not files). The output is piped to 'sort -hr' to sort in descending order using human-readable number comparison, with the largest directories appearing first. Finally, 'head -10' limits output to the top 10 directories.

\section*{Tags}
du, disk-usage, directories, storage, sort, system-administration

\section*{Dependencies}
coreutils


\section*{Arguments}
\begin{enumerate}
  \item \textbf{PATH} (Optional): Directory path to analyze (default: current directory) \\ Default: .
  \item \textbf{COUNT} (Optional): Number of largest directories to display (default: 10) \\ Default: 10
\end{enumerate}

\section*{Examples}
\begin{enumerate}
  \item \texttt{du -sh */ | sort -hr | head -10} - Show top 10 largest subdirectories in current directory
  \item \texttt{du -sh */ | sort -hr} - Show all subdirectories sorted by size without limit
  \item \texttt{du -sh */ | sort -hr | head -5} - Show top 5 largest subdirectories
  \item \texttt{du -sh ./* | sort -hr | head -10} - Include hidden files and directories (those starting with dot)
  \item \texttt{cd /var \&\& du -sh */ | sort -hr | head -10} - Find largest subdirectories in /var
\end{enumerate}

\section*{Output}
\begin{lstlisting}
847G	cache/
512G	data/
256G	backups/
128G	logs/
96G	documents/
64G	downloads/
32G	pictures/
16G	videos/
8.5G	archives/
2.1G	temp/
\end{lstlisting}

\section*{Notes}
\begin{itemize}
  \item The '*/' pattern matches only directories; use '*' to include files as well
  \item The -h flag makes output human-readable (K, M, G, T); use -B for a specific block size
  \item The -s flag summarizes each directory's total size instead of listing contents recursively
  \item Hidden directories (starting with .) are excluded; use './*' or '*' (with shopt) to include them
  \item On systems with many subdirectories, results are calculated quickly since du only traverses each directory once
  \item Symbolic links are not followed by default; add -L flag to follow symlinks
  \item Results may vary if some directories are mounted on different filesystems or are inaccessible due to permissions
\end{itemize}

\section*{Warnings}
\begin{itemize}
  \item Results may be incomplete or misleading if you lack read permissions on some directories
  \item On network filesystems (NFS, SMB), the command may be slower due to network latency
  \item If dealing with millions of directories, the command can consume considerable memory and time
  \item Sparse files may cause inaccurate size reporting on some filesystems
\end{itemize}

\section*{See Also}
\begin{itemize}
  \item find-largest-storage-users
  \item find-large-files-recursive
\end{itemize}

\section*{Status}
reviewed

\section*{Safety}
safe

\section*{Shell}
bash

\section*{Platforms}
linux, gnu-linux, freebsd, openbsd, netbsd

\section*{Created At}
2026-04-21

\section*{Updated At}
2026-04-21

\section{Disk Space Usage Per Directory}
\section*{Display disk space usage for all directories with hierarchical breakdown}
\textbf{Author:} marcos \hfill \textbf{Date:} 2026-04-21

\noindent Shows disk space usage for directories up to 2 levels deep in the current directory tree, sorted by size in descending order with human-readable format, providing a hierarchical breakdown of disk consumption.

\section*{Language}
bash

\section*{Category}
disk-usage

\section*{Command}
\begin{lstlisting}
du -h --max-depth=2 . | sort -hr
\end{lstlisting}

\section*{Explanation}
The command uses 'du -h' to calculate disk usage in human-readable format. The '--max-depth=2' flag limits output to directories up to 2 levels deep from the current directory, avoiding excessive detail while still showing the hierarchy. The '.' specifies the starting directory. The output is piped to 'sort -hr' to sort in descending order by human-readable sizes, making it easy to identify which directories consume the most space at each level.

\section*{Tags}
du, disk-usage, directories, storage, hierarchical, sort, system-administration

\section*{Dependencies}
coreutils


\section*{Arguments}
\begin{enumerate}
  \item \textbf{PATH} (Optional): Root directory to analyze (default: current directory) \\ Default: .
  \item \textbf{DEPTH} (Optional): Maximum directory depth to traverse (default: 2) \\ Default: 2
\end{enumerate}

\section*{Examples}
\begin{enumerate}
  \item \texttt{du -h --max-depth=2 . | sort -hr} - Show disk usage up to 2 levels deep from current directory
  \item \texttt{du -h --max-depth=1 . | sort -hr} - Show only immediate subdirectories (1 level deep)
  \item \texttt{du -h --max-depth=3 /home | sort -hr} - Show disk usage in /home up to 3 levels deep
  \item \texttt{du -h --max-depth=2 . | sort -hr | head -20} - Show top 20 directories by size (2 levels deep)
  \item \texttt{du -ah --max-depth=2 . | sort -hr | head -15} - Include files in the listing, show top 15 items
\end{enumerate}

\section*{Output}
\begin{lstlisting}
1.2T	.
847G	./data
512G	./data/cache
256G	./data/backups
244G	./media
128G	./media/videos
116G	./media/archives
156G	./documents
96G	./documents/projects
60G	./documents/reports
\end{lstlisting}

\section*{Notes}
\begin{itemize}
  \item The --max-depth flag controls recursion depth: 0 shows only the starting directory, 1 shows immediate children, etc.
  \item The -h flag makes output human-readable (K, M, G, T); use -B 1M for sizes in fixed block units
  \item The -s flag is not needed here as du summarizes by default
  \item Hidden directories (those starting with .) are included automatically
  \item Each line shows cumulative size for that directory and all its contents
  \item Symbolic links are not followed by default; add -L to follow symlinks
  \item On filesystems with many nested directories, increasing --max-depth may significantly increase runtime
\end{itemize}

\section*{Warnings}
\begin{itemize}
  \item Results may be incomplete if you lack read permissions on some directories
  \item On network filesystems (NFS, SMB), this command can be slow due to network latency
  \item Very deep directory hierarchies combined with high --max-depth values may cause excessive output
  \item Sparse files may report misleading sizes on some filesystems
\end{itemize}

\section*{See Also}
\begin{itemize}
  \item disk-space-sort-largest-directories
\end{itemize}

\section*{Status}
reviewed

\section*{Safety}
safe

\section*{Shell}
bash

\section*{Platforms}
linux, gnu-linux, freebsd, openbsd, netbsd

\section*{Created At}
2026-04-21

\section*{Updated At}
2026-04-21

\section{Dmesg Errors Pretty}
\section*{Kernel errors/warnings summary with full severity levels and colorized output}
\textbf{Author:} Marcos de Carvalho \hfill \textbf{Date:} 2026-04-21

\noindent Retrieves all kernel messages from the kernel ring buffer, extracts and counts all message severity levels (emerg, alert, crit, err, warn, notice, info, debug), and presents them with a comprehensive summary frontmatter followed by the human-readable colorized logs.

\section*{Language}
bash

\section*{Category}
diagnostics

\section*{Command}
\begin{lstlisting}
dmesg -T -x --level=emerg,alert,crit,err,warn --color=always 2>/dev/null | awk 'BEGIN { print "\033[1m=== KERNEL MESSAGE SUMMARY ===\033[0m" } { lines[NR] = \$0; lvl = \$2; gsub(/\x1B\[[0-9;]*[mK]/, \"\", lvl); gsub(/[: \t]/, \"\", lvl); if(lvl ~ /^(emerg|alert|crit|err|warn|notice|info|debug)\$/) cnt[lvl]++; if(lvl ~ /^\w+\$/) all_cnt[lvl]++ } END { if(NR>0) { print "\033[1m--- ERRORS & WARNINGS ---\033[0m"; for(l in cnt) printf \"  %-7s : %d\n\", l, cnt[l]; print "\033[1m--- ALL MESSAGES (by severity) ---\033[0m"; for(l in all_cnt) printf \"  %-7s : %d\n\", l, all_cnt[l]; print "\033[1m======================================\033[0m\n"; for(i=1; i<=NR; i++) print lines[i] } else print "No kernel messages found for selected levels." }'
\end{lstlisting}

\section*{Explanation}
The one-liner uses 'dmesg -T -x' to extract kernel messages with human-readable timestamps (-T) and decode the facility/level prefixes (-x). We filter for all severity levels (emerg,alert,crit,err,warn,notice,info,debug) via '--level'. The '--color=always' flag preserves standard color coding. The output is piped into 'awk', which does several tasks: (1) strips ANSI escape sequences from the second column to cleanly count occurrences, (2) maintains two counters: one for 'errors/warnings' (emerg,alert,crit,err,warn) and another for 'all messages' (all severity levels), and (3) stores all log lines in memory. The END block prints a formatted summary frontmatter with both the 'errors/warnings' and 'all messages' counts, followed by the original colorized logs.

\section*{Tags}
dmesg, kernel, errors, warnings, logs, troubleshooting, monitoring, summary, awk, diagnostics, severity, log-analysis, linux, system, administration

\section*{Dependencies}
util-linux, gawk


\section*{Arguments}
None

\section*{Examples}
\begin{enumerate}
  \item \texttt{dmesg -T -x --level=emerg,alert,crit,err,warn --color=always 2>/dev/null | awk 'BEGIN \{ print "\textbackslash\{\}033[1m=== KERNEL MESSAGE SUMMARY ===\textbackslash\{\}033[0m" \} \{ lines[NR] = \textbackslash\{\}\$0; lvl = \textbackslash\{\}\$2; gsub(/\textbackslash\{\}x1B\textbackslash\{\}[[0-9;]*[mK]/, \textbackslash\{\}"\textbackslash\{\}", lvl); gsub(/[: \textbackslash\{\}t]/, \textbackslash\{\}"\textbackslash\{\}", lvl); if(lvl \textasciitilde{} /\textasciicircum{}(emerg|alert|crit|err|warn|notice|info|debug)\textbackslash\{\}\$/) cnt[lvl]++; if(lvl \textasciitilde{} /\textasciicircum{}\textbackslash\{\}w+\textbackslash\{\}\$/) all\_cnt[lvl]++ \} END \{ if(NR>0) \{ print "\textbackslash\{\}033[1m--- ERRORS \& WARNINGS ---\textbackslash\{\}033[0m"; for(l in cnt) printf \textbackslash\{\}"  \%-7s : \%d\textbackslash\{\}n\textbackslash\{\}", l, cnt[l]; print "\textbackslash\{\}033[1m--- ALL MESSAGES (by severity) ---\textbackslash\{\}033[0m"; for(l in all\_cnt) printf \textbackslash\{\}"  \%-7s : \%d\textbackslash\{\}n\textbackslash\{\}", l, all\_cnt[l]; print "\textbackslash\{\}033[1m======================================\textbackslash\{\}033[0m\textbackslash\{\}n"; for(i=1; i<=NR; i++) print lines[i] \} else print "No kernel messages found for selected levels." \}'} - Show full severity summary and colorized kernel messages
  \item \texttt{dmesg -T -x --level=emerg,alert,crit,err,warn --color=always 2>/dev/null | awk 'BEGIN \{ print "\textbackslash\{\}033[1m=== KERNEL MESSAGE SUMMARY ===\textbackslash\{\}033[0m" \} \{ lines[NR] = \textbackslash\{\}\$0; lvl = \textbackslash\{\}\$2; gsub(/\textbackslash\{\}x1B\textbackslash\{\}[[0-9;]*[mK]/, \textbackslash\{\}"\textbackslash\{\}", lvl); gsub(/[: \textbackslash\{\}t]/, \textbackslash\{\}"\textbackslash\{\}", lvl); if(lvl \textasciitilde{} /\textasciicircum{}(emerg|alert|crit|err|warn|notice|info|debug)\textbackslash\{\}\$/) cnt[lvl]++; if(lvl \textasciitilde{} /\textasciicircum{}\textbackslash\{\}w+\textbackslash\{\}\$/) all\_cnt[lvl]++ \} END \{ if(NR>0) \{ print "\textbackslash\{\}033[1m--- ERRORS \& WARNINGS ---\textbackslash\{\}033[0m"; for(l in cnt) printf \textbackslash\{\}"  \%-7s : \%d\textbackslash\{\}n\textbackslash\{\}", l, cnt[l]; print "\textbackslash\{\}033[1m--- ALL MESSAGES (by severity) ---\textbackslash\{\}033[0m"; for(l in all\_cnt) printf \textbackslash\{\}"  \%-7s : \%d\textbackslash\{\}n\textbackslash\{\}", l, all\_cnt[l]; print "\textbackslash\{\}033[1m======================================\textbackslash\{\}033[0m\textbackslash\{\}n"; for(i=1; i<=NR; i++) print lines[i] \} else print "No kernel messages found for selected levels." \}' | less -R} - Show the summary and logs with a pager, preserving colors
  \item \texttt{sudo dmesg -T -x --level=emerg,alert,crit,err,warn --color=always 2>/dev/null | awk 'BEGIN \{ print "\textbackslash\{\}033[1m=== KERNEL MESSAGE SUMMARY ===\textbackslash\{\}033[0m" \} \{ lines[NR] = \textbackslash\{\}\$0; lvl = \textbackslash\{\}\$2; gsub(/\textbackslash\{\}x1B\textbackslash\{\}[[0-9;]*[mK]/, \textbackslash\{\}"\textbackslash\{\}", lvl); gsub(/[: \textbackslash\{\}t]/, \textbackslash\{\}"\textbackslash\{\}", lvl); if(lvl \textasciitilde{} /\textasciicircum{}(emerg|alert|crit|err|warn|notice|info|debug)\textbackslash\{\}\$/) cnt[lvl]++; if(lvl \textasciitilde{} /\textasciicircum{}\textbackslash\{\}w+\textbackslash\{\}\$/) all\_cnt[lvl]++ \} END \{ if(NR>0) \{ print "\textbackslash\{\}033[1m--- ERRORS \& WARNINGS ---\textbackslash\{\}033[0m"; for(l in cnt) printf \textbackslash\{\}"  \%-7s : \%d\textbackslash\{\}n\textbackslash\{\}", l, cnt[l]; print "\textbackslash\{\}033[1m--- ALL MESSAGES (by severity) ---\textbackslash\{\}033[0m"; for(l in all\_cnt) printf \textbackslash\{\}"  \%-7s : \%d\textbackslash\{\}n\textbackslash\{\}", l, all\_cnt[l]; print "\textbackslash\{\}033[1m======================================\textbackslash\{\}033[0m\textbackslash\{\}n"; for(i=1; i<=NR; i++) print lines[i] \} else print "No kernel messages found for selected levels." \}'} - Run with sudo to access kernel ring buffer on restricted systems
\end{enumerate}

\section*{Output}
\begin{lstlisting}
\u001b[1m=== KERNEL MESSAGE SUMMARY ===\u001b[0m
\u001b[1m--- ERRORS & WARNINGS ---\u001b[0m
  err     : 2
  warn    : 3
  crit    : 1
\u001b[1m--- ALL MESSAGES (by severity) ---\u001b[0m
  err     : 2
  warn    : 3
  crit    : 1
  notice  : 5
  info    : 12
  debug   : 0
\u001b[1m======================================\u001b[0m

kern  :err   : [Tue Apr 21 10:00:00 2026] sd 0:0:0:0: [sda] Write cache: enabled
kern  :warn  : [Tue Apr 21 10:01:23 2026] WARNING: CPU: 0 PID: 1234
kern  :crit  : [Tue Apr 21 10:02:45 2026] Critical hardware error
kern  :warn  : [Tue Apr 21 10:05:00 2026] ext4-fs: mounted with warnings
kern  :err   : [Tue Apr 21 10:06:12 2026] usb 1-1: device descriptor read/64, error -71
kern  :warn  : [Tue Apr 21 10:08:34 2026] IPv6: ADDRCONF(NETDEV_UP): eth0: link is not ready
kern  :notice: [Tue Apr 21 10:10:00 2026] Kernel log system initialized
kern  :info  : [Tue Apr 21 10:15:00 2026] Network interface eth0 up and running
\end{lstlisting}

\section*{Notes}
\begin{itemize}
  \item Requires read access to the kernel ring buffer, which often requires sudo or being in the 'adm' group on some distributions.
  \item The '-x' (--decode) flag is critical here as it explicitly outputs the severity level on every line, making awk parsing reliable.
  \item ANSI stripping regex in awk ensures color codes don't interfere with severity counting.
  \item This command counts all message severity levels: emerg, alert, crit, err, warn, notice, info, debug.
  \item The 'errors \& warnings' summary includes: emerg, alert, crit, err, warn.
  \item The 'all messages' summary includes all severity levels, providing complete log statistics.
  \item The entire dmesg output is stored in awk's memory (the 'lines' array) before printing, which is perfectly safe for typical volumes but could consume slightly more memory on massively flooded systems.
\end{itemize}

\section*{Warnings}
\begin{itemize}
  \item Running without sufficient privileges (like sudo) will result in a blank output or a permission denied error on most modern secure Linux distributions.
\end{itemize}

\section*{See Also}
\begin{itemize}
  \item find-large-files-recursive
  \item disk-space-sort-largest-directories
\end{itemize}

\section*{Status}
reviewed

\section*{Safety}
safe

\section*{Shell}
bash

\section*{Platforms}
linux

\section*{Created At}
2026-04-21

\section*{Updated At}
2026-04-21

\section{Find 10 Largest Storage Users}
\section*{Find users with largest storage usage}
\textbf{Author:} Marcos de Carvalho \hfill \textbf{Date:} 2026-04-21

\noindent Displays the top 10 users on the system by storage consumption in their home directories, sorted in descending order by disk usage.

\section*{Language}
bash

\section*{Category}
disk-usage

\section*{Command}
\begin{lstlisting}
du -sh /home/* 2>/dev/null | sort -hr | head -10
\end{lstlisting}

\section*{Explanation}
The command uses 'du -sh' to calculate the total disk usage in human-readable format for each home directory. The '2>/dev/null' suppresses permission denied errors for inaccessible directories. The output is piped to 'sort -hr' to sort in descending order (largest first) using human-readable number comparison. Finally, 'head -10' limits the output to the top 10 users.

\section*{Tags}
disk-usage, storage, users, system-administration, du, sort

\section*{Dependencies}
coreutils, findutils


\section*{Arguments}
\begin{enumerate}
  \item \textbf{PATH} (Optional): Home directory path to analyze (default: /home) \\ Default: /home
  \item \textbf{COUNT} (Optional): Number of top users to display (default: 10) \\ Default: 10
\end{enumerate}

\section*{Examples}
\begin{enumerate}
  \item \texttt{du -sh /home/* 2>/dev/null | sort -hr | head -10} - Show top 10 users by storage usage in /home
  \item \texttt{du -sh /home/* 2>/dev/null | sort -hr} - Show all users sorted by storage usage without limit
  \item \texttt{du -sh /home/* 2>/dev/null | sort -hr | head -5} - Show top 5 users by storage usage
  \item \texttt{du -sh /root /home/* 2>/dev/null | sort -hr | head -10} - Include root's home directory in the analysis
\end{enumerate}

\section*{Output}
\begin{lstlisting}
450G	/home/alice
320G	/home/bob
215G	/home/charlie
156G	/home/diana
98G	/home/eve
67G	/home/frank
45G	/home/grace
32G	/home/henry
28G	/home/iris
19G	/home/jack
\end{lstlisting}

\section*{Notes}
\begin{itemize}
  \item The command suppresses permission errors with '2>/dev/null', allowing users to see results without elevated privileges
  \item Human-readable format (-h) displays sizes in K, M, G, etc., making interpretation easier
  \item Sort with -h flag properly handles human-readable numbers (not lexicographic sorting)
  \item Symbolic links in /home are counted separately if they point outside /home
  \item Hidden files and directories within user homes are included in the totals
  \item Results may vary if some home directories are mounted on different filesystems
\end{itemize}

\section*{Warnings}
\begin{itemize}
  \item Results may be incomplete if you lack read permissions on certain home directories
  \item On systems with many users, this command may take considerable time to complete
  \item NFS-mounted or slow storage may cause noticeable delays
\end{itemize}

\section*{See Also}
\begin{itemize}
  \item find-large-files-recursive
  \item disk-space-usage-per-directory
\end{itemize}

\section*{Status}
reviewed

\section*{Safety}
safe

\section*{Shell}
bash

\section*{Platforms}
linux, gnu-linux, freebsd

\section*{Created At}
2026-04-21

\section*{Updated At}
2026-04-21

\section{Find Large Files Recursive}
\section*{Find large files recursively in a directory tree}
\textbf{Author:} marcos \hfill \textbf{Date:} 2026-04-21

\noindent Recursively finds all regular files larger than 100 MB in the directory tree, displays them with human-readable sizes, and sorts by size in descending order.

\section*{Language}
bash

\section*{Category}
filesystem

\section*{Command}
\begin{lstlisting}
find . -type f -size +100M -exec ls -lh {} + | awk '{print $5, $9}' | sort -hr
\end{lstlisting}

\section*{Explanation}
The command uses 'find' to recursively traverse the directory tree, filter for regular files (-type f), and select those exceeding 100 MB (-size +100M). The '-exec ls -lh \{\} +' displays detailed information for each file in a batch operation, which is more efficient than using '\textbackslash\{\};'. The output is piped to 'awk' to extract the file size and path, then 'sort -hr' sorts the results by human-readable sizes in descending order, showing the largest files first.

\section*{Tags}
find, filesystem, disk-usage, large-files, recursive, storage

\section*{Dependencies}
findutils, coreutils


\section*{Arguments}
\begin{enumerate}
  \item \textbf{PATH} (Optional): Root directory to start the recursive search (default: current directory) \\ Default: .
  \item \textbf{SIZE} (Optional): Minimum file size threshold (default: 100M, supports K, M, G suffixes) \\ Default: 100M
\end{enumerate}

\section*{Examples}
\begin{enumerate}
  \item \texttt{find . -type f -size +100M -exec ls -lh \{\} + | awk '\{print \$5, \$9\}' | sort -hr} - Find files larger than 100 MB in current directory and subdirectories
  \item \texttt{find /var -type f -size +500M -exec ls -lh \{\} + | awk '\{print \$5, \$9\}' | sort -hr} - Find large files in /var directory
  \item \texttt{find . -type f -size +1G -exec ls -lh \{\} + | awk '\{print \$5, \$9\}' | sort -hr} - Find files larger than 1 GB
  \item \texttt{find . -type f -size +50M -exec ls -lh \{\} + | awk '\{print \$5, \$9\}' | sort -hr | head -20} - Show top 20 largest files over 50 MB
\end{enumerate}

\section*{Output}
\begin{lstlisting}
1.5G ./videos/archive.tar.gz
987M ./backups/database.sql.bz2
756M ./datasets/training-data.bin
543M ./cache/node_modules.tar
387M ./downloads/ubuntu-22.04-iso
256M ./logs/app-2026.log
198M ./tmp/large-temp-file
\end{lstlisting}

\section*{Notes}
\begin{itemize}
  \item The -size option supports various units: c (bytes), k (kilobytes), M (megabytes), G (gigabytes), and b (512-byte blocks)
  \item By default, -size rounds up (e.g., +100M matches files larger than 100 MiB)
  \item Use -size -100M to find files smaller than 100 MB
  \item The command may take considerable time on large directory trees
  \item Using '\{\}  +' instead of '\{\}  \textbackslash\{\};' is more efficient as it batches the command executions
  \item Results may be incomplete in directories where you lack read permissions
\end{itemize}

\section*{Warnings}
\begin{itemize}
  \item On network filesystems (NFS, SMB), this command may be significantly slower
  \item Does not follow symbolic links by default; add -L flag if needed
  \item Very large directory trees may cause the command to run for extended periods
  \item Some filesystems may not report accurate sizes for sparse files
\end{itemize}

\section*{See Also}
\begin{itemize}
  \item find-largest-storage-users
  \item disk-space-sort-largest-directories
\end{itemize}

\section*{Status}
reviewed

\section*{Safety}
safe

\section*{Shell}
bash

\section*{Platforms}
linux, gnu-linux, freebsd, openbsd, netbsd

\section*{Created At}
2026-04-21

\section*{Updated At}
2026-04-21

\section{Journalctl Errors Today}
\section*{Show system executables with errors/failures from today's journal}
\textbf{Author:} marcos \hfill \textbf{Date:} 2026-04-22

\noindent Lists system executables that have generated failure, error, or fatal messages in today's journal logs, sorted by frequency.

\section*{Language}
bash

\section*{Category}
system-monitoring

\section*{Command}
\begin{lstlisting}
echo "=== Errors/Failures in Today's Journal ==="; journalctl --no-pager --since today --grep 'fail|error|fatal' --output json | jq -r '._EXE' | sort | uniq -c | sort -nr | awk '{printf "%4d  %s\n", $1, $2}'
\end{lstlisting}

\section*{Explanation}
The command retrieves all journal entries from today (since midnight), filters for messages containing the keywords 'fail', 'error', or 'fatal' (case-insensitive regex), outputs the results in JSON format for parsing, extracts the executable name (\_EXE field) from each log entry, counts occurrences of each executable, sorts them numerically in descending order (most frequent first), and formats the output with a header and aligned columns for readability.

\section*{Tags}
journalctl, systemd, logs, monitoring, troubleshooting, errors, failure, system-administration, diagnostics, debugging, logging, bash, jq, text-processing, process-management

\section*{Dependencies}
systemd, jq


\section*{Arguments}
None

\section*{Examples}
\begin{enumerate}
  \item \texttt{journalctl-errors-today} - Show errors/failures from today's journal
  \item \texttt{journalctl --no-pager --since yesterday --grep 'fail|error|fatal' --output json | jq -r '.\_EXE' | sort | uniq -c | sort -nr} - Show errors/failures from yesterday's journal (raw output)
\end{enumerate}

\section*{Output}
\begin{lstlisting}
=== Errors/Failures in Today's Journal ===
   5  /usr/bin/systemd
   3  /usr/bin/foo-daemon
   1  /usr/lib/bar-service
\end{lstlisting}

\section*{Notes}
\begin{itemize}
  \item Requires journalctl (systemd) and jq to be installed
  \item Only shows entries from the current day (midnight to now)
  \item Requires appropriate permissions to read journal logs
  \item The regex pattern 'fail|error|fatal' matches any of these words in log messages
  \item Output is sorted by count (highest first) for easy identification of frequent issues
\end{itemize}

\section*{Warnings}
None

\section*{See Also}
\begin{itemize}
  \item dmesg-errors-pretty
  \item user-activity-and-quota-report
\end{itemize}

\section*{Status}
reviewed

\section*{Safety}
safe

\section*{Shell}
bash

\section*{Platforms}
linux, gnu-linux

\section*{Created At}
2026-04-22

\section*{Updated At}
2026-04-22

\section{Nonhuman User Activity And Quota Report}
\section*{Comprehensive system/service non-human user report with disk quota, usage, and activity}
\textbf{Author:} Marcos de Carvalho \hfill \textbf{Date:} 2026-04-21

\noindent Generates a comprehensive system audit report for all non-human (system/service) users (requires sudo) showing username, disk quota limit, home directory path, actual disk usage, account creation date, file count, last login timestamp, and last executed command. All information is formatted in a readable table with elevated privileges for complete data access.

\section*{Language}
bash

\section*{Category}
system-administration

\section*{Command}
\begin{lstlisting}
sudo bash -c "printf 'USER\tQUOTA\tHOME\tDISK_USAGE\tCREATED\tFILES\tLAST_LOGIN\tLAST_CMD\n'; getent passwd | awk -F: '\$3 < 1000 || \$1 ~ /^nobody$/ { print \$1,\$6 }' | while read u h; do q=\$(quota -u \"\$u\" 2>/dev/null | tail -1 | awk '{print \$2}'); [ -z \"\$q\" ] && q=none; d=\$(du -sh \"\$h\" 2>/dev/null | cut -f1); c=\$(stat -c %y \"\$h\" 2>/dev/null | cut -d' ' -f1 || echo N/A); f=\$(find \"\$h\" -type f 2>/dev/null | wc -l); l=\$(lastlog -u \"\$u\" 2>/dev/null | tail -1 | awk '{print \$5,\$6,\$7}'); [ -z \"\$l\" ] && l=Never; cmd=\$(tail -1 \"\$h\"/.bash_history 2>/dev/null | head -c 30); [ -z \"\$cmd\" ] && cmd=N/A; printf '%s\t%s\t%s\t%s\t%s\t%s\t%s\t%s\n' \"\$u\" \"\$q\" \"\$h\" \"\$d\" \"\$c\" \"\$f\" \"\$l\" \"\$cmd\"; done" | column -t -s $'\t'
\end{lstlisting}

\section*{Explanation}
The command uses 'sudo bash -c' to execute the entire script with elevated privileges, ensuring access to all user data. It retrieves users with UID < 1000 (or 'nobody') via getent, then for each user gathers: disk quota from 'quota -u', home disk usage from 'du -sh', account creation date via 'stat', file count using 'find', last login from 'lslogins', and last command from .bash\_history. The 'bash -c' wrapper allows proper variable expansion and quoting within the sudo context. Output is joined by tabs and formatted by 'column -t' for readability, avoiding issues with spaces in values.

\section*{Tags}
user-management, disk-usage, system-administration, activity, report, quota, audit, system-users, service-accounts

\section*{Dependencies}
coreutils, util-linux, shadow-utils, findutils


\section*{Arguments}
\begin{enumerate}
  \item \textbf{UID\_THRESHOLD} (Optional): Maximum UID to consider a system user (default: 1000 for non-human users) \\ Default: 1000
\end{enumerate}

\section*{Examples}
\begin{enumerate}
  \item \texttt{sudo bash -c "printf 'USER\textbackslash\{\}tQUOTA\textbackslash\{\}tHOME\textbackslash\{\}tDISK\_USAGE\textbackslash\{\}tCREATED\textbackslash\{\}tFILES\textbackslash\{\}tLAST\_LOGIN\textbackslash\{\}tLAST\_CMD\textbackslash\{\}n'; getent passwd | awk -F: '\textbackslash\{\}\$3 < 1000 || \textbackslash\{\}\$1 \textasciitilde{} /\textasciicircum{}nobody\$/ \{ print \textbackslash\{\}\$1,\textbackslash\{\}\$6 \}' | while read u h; do q=\textbackslash\{\}\$(quota -u \textbackslash\{\}"\textbackslash\{\}\$u\textbackslash\{\}" 2>/dev/null | tail -1 | awk '\{print \textbackslash\{\}\$2\}'); [ -z \textbackslash\{\}"\textbackslash\{\}\$q\textbackslash\{\}" ] \&\& q=none; d=\textbackslash\{\}\$(du -sh \textbackslash\{\}"\textbackslash\{\}\$h\textbackslash\{\}" 2>/dev/null | cut -f1); c=\textbackslash\{\}\$(stat -c \%y \textbackslash\{\}"\textbackslash\{\}\$h\textbackslash\{\}" 2>/dev/null | cut -d' ' -f1 || echo N/A); f=\textbackslash\{\}\$(find \textbackslash\{\}"\textbackslash\{\}\$h\textbackslash\{\}" -type f 2>/dev/null | wc -l); l=\textbackslash\{\}\$(lastlog -u \textbackslash\{\}"\textbackslash\{\}\$u\textbackslash\{\}" 2>/dev/null | tail -1 | awk '\{print \textbackslash\{\}\$5,\textbackslash\{\}\$6,\textbackslash\{\}\$7\}'); [ -z \textbackslash\{\}"\textbackslash\{\}\$l\textbackslash\{\}" ] \&\& l=Never; cmd=\textbackslash\{\}\$(tail -1 \textbackslash\{\}"\textbackslash\{\}\$h\textbackslash\{\}"/.bash\_history 2>/dev/null | head -c 30); [ -z \textbackslash\{\}"\textbackslash\{\}\$cmd\textbackslash\{\}" ] \&\& cmd=N/A; printf '\%s\textbackslash\{\}t\%s\textbackslash\{\}t\%s\textbackslash\{\}t\%s\textbackslash\{\}t\%s\textbackslash\{\}t\%s\textbackslash\{\}t\%s\textbackslash\{\}t\%s\textbackslash\{\}n' \textbackslash\{\}"\textbackslash\{\}\$u\textbackslash\{\}" \textbackslash\{\}"\textbackslash\{\}\$q\textbackslash\{\}" \textbackslash\{\}"\textbackslash\{\}\$h\textbackslash\{\}" \textbackslash\{\}"\textbackslash\{\}\$d\textbackslash\{\}" \textbackslash\{\}"\textbackslash\{\}\$c\textbackslash\{\}" \textbackslash\{\}"\textbackslash\{\}\$f\textbackslash\{\}" \textbackslash\{\}"\textbackslash\{\}\$l\textbackslash\{\}" \textbackslash\{\}"\textbackslash\{\}\$cmd\textbackslash\{\}"; done" | column -t -s \$'\textbackslash\{\}t'} - Generate complete report for all non-human (system) users with all metrics (requires sudo)
  \item \texttt{sudo bash -c "printf 'USER\textbackslash\{\}tQUOTA\textbackslash\{\}tDISK\_USAGE\textbackslash\{\}tCREATED\textbackslash\{\}tFILES\textbackslash\{\}tLAST\_LOGIN\textbackslash\{\}n'; getent passwd | awk -F: '\textbackslash\{\}\$3 < 1000 || \textbackslash\{\}\$1 \textasciitilde{} /\textasciicircum{}nobody\$/ \{ print \textbackslash\{\}\$1,\textbackslash\{\}\$6 \}' | while read u h; do q=\textbackslash\{\}\$(quota -u \textbackslash\{\}"\textbackslash\{\}\$u\textbackslash\{\}" 2>/dev/null | tail -1 | awk '\{print \textbackslash\{\}\$2\}'); [ -z \textbackslash\{\}"\textbackslash\{\}\$q\textbackslash\{\}" ] \&\& q=none; d=\textbackslash\{\}\$(du -sh \textbackslash\{\}"\textbackslash\{\}\$h\textbackslash\{\}" 2>/dev/null | cut -f1); c=\textbackslash\{\}\$(stat -c \%y \textbackslash\{\}"\textbackslash\{\}\$h\textbackslash\{\}" 2>/dev/null | cut -d' ' -f1 || echo N/A); f=\textbackslash\{\}\$(find \textbackslash\{\}"\textbackslash\{\}\$h\textbackslash\{\}" -type f 2>/dev/null | wc -l); l=\textbackslash\{\}\$(lastlog -u \textbackslash\{\}"\textbackslash\{\}\$u\textbackslash\{\}" 2>/dev/null | tail -1 | awk '\{print \textbackslash\{\}\$5,\textbackslash\{\}\$6,\textbackslash\{\}\$7\}'); [ -z \textbackslash\{\}"\textbackslash\{\}\$l\textbackslash\{\}" ] \&\& l=Never; printf '\%s\textbackslash\{\}t\%s\textbackslash\{\}t\%s\textbackslash\{\}t\%s\textbackslash\{\}t\%s\textbackslash\{\}t\%s\textbackslash\{\}n' \textbackslash\{\}"\textbackslash\{\}\$u\textbackslash\{\}" \textbackslash\{\}"\textbackslash\{\}\$q\textbackslash\{\}" \textbackslash\{\}"\textbackslash\{\}\$d\textbackslash\{\}" \textbackslash\{\}"\textbackslash\{\}\$c\textbackslash\{\}" \textbackslash\{\}"\textbackslash\{\}\$f\textbackslash\{\}" \textbackslash\{\}"\textbackslash\{\}\$l\textbackslash\{\}"; done | sort -t \$'\textbackslash\{\}t' -k3 -hr" | column -t -s \$'\textbackslash\{\}t'} - Sort by disk usage descending, exclude home path and last command
  \item \texttt{sudo bash -c "printf 'USER\textbackslash\{\}tQUOTA\textbackslash\{\}tDISK\_USAGE\textbackslash\{\}tFILES\textbackslash\{\}n'; getent passwd | awk -F: '\textbackslash\{\}\$3 < 1000 || \textbackslash\{\}\$1 \textasciitilde{} /\textasciicircum{}nobody\$/ \{ print \textbackslash\{\}\$1,\textbackslash\{\}\$6 \}' | while read u h; do q=\textbackslash\{\}\$(quota -u \textbackslash\{\}"\textbackslash\{\}\$u\textbackslash\{\}" 2>/dev/null | tail -1 | awk '\{print \textbackslash\{\}\$2\}'); [ -z \textbackslash\{\}"\textbackslash\{\}\$q\textbackslash\{\}" ] \&\& q=none; d=\textbackslash\{\}\$(du -sh \textbackslash\{\}"\textbackslash\{\}\$h\textbackslash\{\}" 2>/dev/null | cut -f1); f=\textbackslash\{\}\$(find \textbackslash\{\}"\textbackslash\{\}\$h\textbackslash\{\}" -type f 2>/dev/null | wc -l); printf '\%s\textbackslash\{\}t\%s\textbackslash\{\}t\%s\textbackslash\{\}t\%s\textbackslash\{\}n' \textbackslash\{\}"\textbackslash\{\}\$u\textbackslash\{\}" \textbackslash\{\}"\textbackslash\{\}\$q\textbackslash\{\}" \textbackslash\{\}"\textbackslash\{\}\$d\textbackslash\{\}" \textbackslash\{\}"\textbackslash\{\}\$f\textbackslash\{\}"; done" | column -t -s \$'\textbackslash\{\}t'} - Simplified report showing only quota, disk usage, and file count
\end{enumerate}

\section*{Output}
\begin{lstlisting}
USER        QUOTA HOME            DISK_USAGE CREATED    FILES  LAST_LOGIN      LAST_CMD
root        none  /root           1.5G       2023-01-15 15340  May  20 13:45   apt update
daemon      none  /usr/sbin       32K        2023-01-15 12     Never           N/A
bin         none  /bin            0          2023-01-15 0      Never           N/A
www-data    none  /var/www        4.2G       2023-05-10 45678  Never           N/A
nobody      none  /nonexistent    0          2023-01-15 0      Never           N/A
\end{lstlisting}

\section*{Notes}
\begin{itemize}
  \item This command is designed to run with sudo for complete information; running without sudo will show incomplete results
  \item The 'sudo bash -c' wrapper preserves the entire command context and allows proper variable expansion
  \item Disk quota requires quota system to be enabled on the filesystem; shows 'none' if unavailable
  \item File count calculation uses 'find' which can be slow on directories with millions of files
  \item Account creation date uses home directory birth time (\%w) if available, falling back to modification time (\%y)
  \item Last command is truncated to 30 characters; extracted from .bash\_history. Note: Bash only flushes history to disk upon session exit or when 'history -a' is run, so this shows the last saved command, not necessarily the active session's most recent command.
  \item Users with UIDs >= 1000 are standard human accounts and are excluded by default. Note that system/service accounts typically lack valid login shells or history files, making LAST\_CMD frequently N/A and LAST\_LOGIN Never.
  \item The command uses getent instead of /etc/passwd for better compatibility with different user databases (LDAP, NIS, etc.)
\end{itemize}

\section*{Warnings}
\begin{itemize}
  \item Requires sudo access to function properly; user must have sudo privileges without password prompt configured in sudoers for unattended execution
  \item Large file counts or extensive directory trees will cause significant performance degradation
  \item Running with find on large home directories can consume considerable IO and CPU resources
  \item The lslogins command retrieves last login information depending on wtmp/btmp logs; accuracy varies by distribution logging policies
  \item Last command information depends on .bash\_history being available and enabled; other shells use different history files
  \item Network-mounted home directories (NFS, SMB) will significantly increase command execution time
  \item On systems with many users, overall execution time can be substantial; consider running during off-peak hours
\end{itemize}

\section*{See Also}
\begin{itemize}
  \item find-largest-storage-users
  \item du-sort-largest-directories
  \item disk-space-usage-per-directory
  \item find-large-files-recursive
\end{itemize}

\section*{Status}
reviewed

\section*{Safety}
caution

\section*{Shell}
bash

\section*{Platforms}
linux, gnu-linux

\section*{Created At}
2026-04-21

\section*{Updated At}
2026-04-21

\section{Replace String In Files}
\section*{Replace string in multiple files using sed}
\textbf{Author:} Marcos de Carvalho \hfill \textbf{Date:} 2026-04-21

\noindent Replaces all occurrences of a string with a replacement string in multiple files or file patterns. Pure oneliner that works directly pasted with parameters, or in an alias.

\section*{Language}
bash

\section*{Category}
text-processing

\section*{Command}
\begin{lstlisting}
bash -c 'old=$1; new=$2; shift 2; find "$@" -type f -exec sed -i "s#$old#$new#g" {} +' _
\end{lstlisting}

\section*{Explanation}
This oneliner uses a subshell (`bash -c`) to parse positional parameters. The first argument is the string to replace (`\$1`), the second is the replacement (`\$2`), and `shift 2` leaves the remaining arguments as the files or patterns to search in (`\$@`). The `find` command locates all matching files and processes them with `sed` using the `-exec \{\} +` syntax, which efficiently handles multiple files in batches. The `\_` at the end acts as `\$0` (the script name) so `\$1` maps to the first actual argument. This approach handles an arbitrary amount of files and supports wildcard file patterns without hardcoding values into the command itself.

\section*{Tags}
sed, text-processing, string-replacement, file-editing, search-replace, batch-replace

\section*{Dependencies}
sed, coreutils


\section*{Arguments}
\begin{enumerate}
  \item \textbf{SEARCH} (Required): The string to search for (first parameter passed after the command).
  \item \textbf{REPLACE} (Required): The replacement string (second parameter passed after the command).
  \item \textbf{FILES} (Required): One or more file paths or patterns (e.g., file1.txt, *.txt, /path/*.js).
\end{enumerate}

\section*{Examples}
\begin{enumerate}
  \item \texttt{bash -c 'old=\$1; new=\$2; shift 2; find "\$@" -type f -exec sed -i "s\#\$old\#\$new\#g" \{\} +' \_ 'old-domain.com' 'new-domain.com' config.json settings.json} - Paste directly: replace 'old-domain.com' with 'new-domain.com' in config.json and settings.json
  \item \texttt{bash -c 'old=\$1; new=\$2; shift 2; find "\$@" -type f -exec sed -i "s\#\$old\#\$new\#g" \{\} +' \_ '/var/old/path' '/var/new/path' nginx.conf} - Replace a path containing slashes in a configuration file
  \item \texttt{bash -c 'old=\$1; new=\$2; shift 2; find "\$@" -type f -exec sed -i "s\#\$old\#\$new\#g" \{\} +' \_ 'old-text' 'new-text' *.txt} - Replace in all .txt files in current directory using a wildcard
  \item \texttt{alias repl="bash -c 'old=\textbackslash\{\}\$1; new=\textbackslash\{\}\$2; shift 2; find \textbackslash\{\}"\textbackslash\{\}\$@\textbackslash\{\}" -type f -exec sed -i \textbackslash\{\}"s\#\textbackslash\{\}\$old\#\textbackslash\{\}\$new\#g\textbackslash\{\}" \{\} +' \_"} - Create an alias named 'repl' to easily perform string replacements
  \item \texttt{repl 'deprecated\_method' 'new\_method' src/**/*.py} - Use the alias to replace in all Python files in src directory and subdirectories
  \item \texttt{repl 'old-api' 'new-api' config/*.json docs/*.md} - Use the alias to replace in multiple file types across different directories
\end{enumerate}

\section*{Output}
\begin{lstlisting}
(No output on success - files are modified in place)
\end{lstlisting}

\section*{Notes}
\begin{itemize}
  \item Pure oneliner: uses bash -c to accept arbitrary parameters after the command string
  \item Usage: Paste the command and append your search string, replacement string, and file patterns
  \item Can be easily turned into a powerful alias in your .bashrc or .bash\_profile
  \item The -i flag modifies files in place; make backups first with cp file.txt file.txt.bak if concerned
  \item Using \# as delimiter prevents issues when search/replace strings contain forward slashes
  \item The g flag replaces all occurrences per line; remove it to replace only first occurrence
  \item The -exec \{\} + syntax bundles multiple files to sed for efficiency
  \item Files with spaces or special characters are handled correctly by find
  \item For case-insensitive replacement, modify the command to include the i flag: sed -i "s\#\$old\#\$new\#gi"
\end{itemize}

\section*{Warnings}
\begin{itemize}
  \item The -i flag modifies files directly without confirmation - always test on copies first
  \item If search string is empty, sed will replace every character with the replacement string
  \item The command will silently fail to modify files if they lack write permissions
  \item Very large files or many files may cause performance issues
  \item Replacement is literal - regex special characters will be treated as patterns (not escaped)
  \item Be careful with wildcard patterns - they may match more files than intended
  \item The find command only processes regular files, not directories
  \item Special characters in search/replace strings must be properly quoted or escaped
\end{itemize}

\section*{See Also}
\begin{itemize}
  \item find-large-files-recursive
  \item disk-space-usage-per-directory
\end{itemize}

\section*{Status}
reviewed

\section*{Safety}
caution

\section*{Shell}
bash

\section*{Platforms}
linux, gnu-linux, freebsd, openbsd, netbsd

\section*{Created At}
2026-04-21

\section*{Updated At}
2026-04-21

\section{System Memory Report Top10}
\section*{System memory report with top 10 consuming processes}
\textbf{Author:} Marcos de Carvalho \hfill \textbf{Date:} 2026-04-22

\noindent Displays system memory usage in human-readable format and lists the top 10 memory-consuming processes by percentage of memory used and resident set size in MB.

\section*{Language}
bash

\section*{Category}
monitoring

\section*{Command}
\begin{lstlisting}
echo -e "Memory Report:\n$(free -h)\n\nTop 10 Memory Consuming Processes (%MEM and RSS in MB):\n$(ps -eo pid,comm,%mem,rss --sort=-%mem | awk 'NR==1 {print $0 " RSS(MB)"}; NR>1 {printf "%s %s %s %.2f\n", $1, $2, $3, $4/1024}' | head -n 11)"
\end{lstlisting}

\section*{Explanation}
The command uses 'free -h' to show memory statistics in human-readable units (e.g., MB, GB). It then uses 'ps' to list all processes sorted by memory usage (\%MEM) in descending order, showing the top 10 processes including their RSS (Resident Set Size) converted to MB. The awk command formats the output to display PID, COMMAND, \%MEM, and RSS(MB) with proper headers. The output is formatted with clear section headers and spacing for readability.

\section*{Tags}
memory, system-monitoring, top-processes, ram, ps, free

\section*{Dependencies}
bash, coreutils, procps-ng


\section*{Arguments}
None

\section*{Examples}
\begin{enumerate}
  \item \texttt{echo -e "Memory Report:\textbackslash\{\}n\$(free -h)\textbackslash\{\}n\textbackslash\{\}nTop 10 Memory Consuming Processes (\%MEM and RSS in MB):\textbackslash\{\}n\$(ps -eo pid,comm,\%mem,rss --sort=-\%mem | awk 'NR==1 \{print \$0 " RSS(MB)"\}; NR>1 \{printf "\%s \%s \%s \%.2f\textbackslash\{\}n", \$1, \$2, \$3, \$4/1024\}' | head -n 11)"} - Run the memory report oneliner directly in the shell
  \item \texttt{alias memreport='echo -e "Memory Report:\textbackslash\{\}n\$(free -h)\textbackslash\{\}n\textbackslash\{\}nTop 10 Memory Consuming Processes (\%MEM and RSS in MB):\textbackslash\{\}n\$(ps -eo pid,comm,\%mem,rss --sort=-\%mem | awk 'NR==1 \{print \$0 " RSS(MB)"\}; NR>1 \{printf "\%s \%s \%s \%.2f\textbackslash\{\}n", \$1, \$2, \$3, \$4/1024\}' | head -n 11)"' \&\& memreport} - Define as an alias and then execute it
\end{enumerate}

\section*{Output}
\begin{lstlisting}
Memory Report:
              total        used        free      shared  buff/cache   available
Mem:           7.7G        2.3G        3.1G        256M        2.3G        5.0G
Swap:          2.0G          0B        2.0G

Top 10 Memory Consuming Processes (%MEM and RSS in MB):
  PID COMMAND         %MEM  RSS(MB)
 1234 firefox         15.2  1175.04
 5678 chrome          12.8   991.42
 9012 code             8.5   657.92
 3456 thunderbird      6.2   479.74
 7890 slack            4.1   317.12
 1122 python3          3.8   293.76
 3344 java             3.5   270.85
 5566 docker           2.9   224.26
 7788 Xorg             2.5   193.28
 9900 gnome-shell      2.2   170.24
\end{lstlisting}

\section*{Notes}
\begin{itemize}
  \item The \%MEM column shows the percentage of available physical memory used by the process.
  \item The RSS(MB) column shows the Resident Set Size in megabytes (non-swapped physical memory used).
  \item The header line is included in the top 10 output for clarity.
  \item The awk command converts RSS from KB to MB by dividing by 1024 and formats to 2 decimal places.
\end{itemize}

\section*{Warnings}
None

\section*{See Also}
\begin{itemize}
  \item disk-usage-summary
  \item cpu-top-processes
\end{itemize}

\section*{Status}
reviewed

\section*{Safety}
safe

\section*{Shell}
bash

\section*{Platforms}
linux, gnu-linux, freebsd, openbsd, netbsd

\section*{Created At}
2026-04-22

\section*{Updated At}
2026-04-22

\section{User Activity And Quota Report}
\section*{Comprehensive user system report with disk quota, usage, and activity}
\textbf{Author:} Marcos de Carvalho \hfill \textbf{Date:} 2026-04-21

\noindent Generates a comprehensive system audit report for all human users (requires sudo) showing username, disk quota limit, home directory path, actual disk usage, account creation date, file count, last login timestamp, and last executed command. All information is formatted in a readable table with elevated privileges for complete data access.

\section*{Language}
bash

\section*{Category}
system-administration

\section*{Command}
\begin{lstlisting}
sudo bash -c "printf 'USER\tQUOTA\tHOME\tDISK_USAGE\tCREATED\tFILES\tLAST_LOGIN\tLAST_CMD\n'; getent passwd | awk -F: '\$3 >= 1000 { if (\$1 ~ /^nobody$/) next; print \$1,\$6 }' | while read u h; do q=\$(quota -u \"\$u\" 2>/dev/null | tail -1 | awk '{print \$2}'); [ -z \"\$q\" ] && q=none; d=\$(du -sh \"\$h\" 2>/dev/null | cut -f1); c=\$(stat -c %y \"\$h\" 2>/dev/null | cut -d' ' -f1 || echo N/A); f=\$(find \"\$h\" -type f 2>/dev/null | wc -l); l=\$(lastlog -u \"\$u\" 2>/dev/null | tail -1 | awk '{print \$5,\$6,\$7}'); [ -z \"\$l\" ] && l=Never; cmd=\$(tail -1 \"\$h\"/.bash_history 2>/dev/null | head -c 30); [ -z \"\$cmd\" ] && cmd=N/A; printf '%s\t%s\t%s\t%s\t%s\t%s\t%s\t%s\n' \"\$u\" \"\$q\" \"\$h\" \"\$d\" \"\$c\" \"\$f\" \"\$l\" \"\$cmd\"; done" | column -t -s $'\t'
\end{lstlisting}

\section*{Explanation}
The command uses 'sudo bash -c' to execute the entire script with elevated privileges, ensuring access to all user data. It retrieves users with UID >= 1000 via getent, then for each user gathers: disk quota from 'quota -u', home disk usage from 'du -sh', account creation date via 'stat', file count using 'find', last login from 'lslogins', and last command from .bash\_history. The 'bash -c' wrapper allows proper variable expansion and quoting within the sudo context. Output is joined by tabs and formatted by 'column -t' for readability, avoiding issues with spaces in values.

\section*{Tags}
user-management, disk-usage, system-administration, activity, report, quota, audit, users

\section*{Dependencies}
coreutils, util-linux, shadow-utils, findutils


\section*{Arguments}
\begin{enumerate}
  \item \textbf{UID\_THRESHOLD} (Optional): Minimum UID to consider a user (default: 1000 for human users only) \\ Default: 1000
\end{enumerate}

\section*{Examples}
\begin{enumerate}
  \item \texttt{sudo bash -c "printf 'USER\textbackslash\{\}tQUOTA\textbackslash\{\}tHOME\textbackslash\{\}tDISK\_USAGE\textbackslash\{\}tCREATED\textbackslash\{\}tFILES\textbackslash\{\}tLAST\_LOGIN\textbackslash\{\}tLAST\_CMD\textbackslash\{\}n'; getent passwd | awk -F: '\textbackslash\{\}\$3 >= 1000 \{ if (\textbackslash\{\}\$1 \textasciitilde{} /\textasciicircum{}nobody\$/) next; print \textbackslash\{\}\$1,\textbackslash\{\}\$6 \}' | while read u h; do q=\textbackslash\{\}\$(quota -u \textbackslash\{\}"\textbackslash\{\}\$u\textbackslash\{\}" 2>/dev/null | tail -1 | awk '\{print \textbackslash\{\}\$2\}'); [ -z \textbackslash\{\}"\textbackslash\{\}\$q\textbackslash\{\}" ] \&\& q=none; d=\textbackslash\{\}\$(du -sh \textbackslash\{\}"\textbackslash\{\}\$h\textbackslash\{\}" 2>/dev/null | cut -f1); c=\textbackslash\{\}\$(stat -c \%y \textbackslash\{\}"\textbackslash\{\}\$h\textbackslash\{\}" 2>/dev/null | cut -d' ' -f1 || echo N/A); f=\textbackslash\{\}\$(find \textbackslash\{\}"\textbackslash\{\}\$h\textbackslash\{\}" -type f 2>/dev/null | wc -l); l=\textbackslash\{\}\$(lastlog -u \textbackslash\{\}"\textbackslash\{\}\$u\textbackslash\{\}" 2>/dev/null | tail -1 | awk '\{print \textbackslash\{\}\$5,\textbackslash\{\}\$6,\textbackslash\{\}\$7\}'); [ -z \textbackslash\{\}"\textbackslash\{\}\$l\textbackslash\{\}" ] \&\& l=Never; cmd=\textbackslash\{\}\$(tail -1 \textbackslash\{\}"\textbackslash\{\}\$h\textbackslash\{\}"/.bash\_history 2>/dev/null | head -c 30); [ -z \textbackslash\{\}"\textbackslash\{\}\$cmd\textbackslash\{\}" ] \&\& cmd=N/A; printf '\%s\textbackslash\{\}t\%s\textbackslash\{\}t\%s\textbackslash\{\}t\%s\textbackslash\{\}t\%s\textbackslash\{\}t\%s\textbackslash\{\}t\%s\textbackslash\{\}t\%s\textbackslash\{\}n' \textbackslash\{\}"\textbackslash\{\}\$u\textbackslash\{\}" \textbackslash\{\}"\textbackslash\{\}\$q\textbackslash\{\}" \textbackslash\{\}"\textbackslash\{\}\$h\textbackslash\{\}" \textbackslash\{\}"\textbackslash\{\}\$d\textbackslash\{\}" \textbackslash\{\}"\textbackslash\{\}\$c\textbackslash\{\}" \textbackslash\{\}"\textbackslash\{\}\$f\textbackslash\{\}" \textbackslash\{\}"\textbackslash\{\}\$l\textbackslash\{\}" \textbackslash\{\}"\textbackslash\{\}\$cmd\textbackslash\{\}"; done" | column -t -s \$'\textbackslash\{\}t'} - Generate complete report for all human users with all metrics (requires sudo)
  \item \texttt{sudo bash -c "printf 'USER\textbackslash\{\}tQUOTA\textbackslash\{\}tDISK\_USAGE\textbackslash\{\}tCREATED\textbackslash\{\}tFILES\textbackslash\{\}tLAST\_LOGIN\textbackslash\{\}n'; getent passwd | awk -F: '\textbackslash\{\}\$3 >= 1000 \{ if (\textbackslash\{\}\$1 \textasciitilde{} /\textasciicircum{}nobody\$/) next; print \textbackslash\{\}\$1,\textbackslash\{\}\$6 \}' | while read u h; do q=\textbackslash\{\}\$(quota -u \textbackslash\{\}"\textbackslash\{\}\$u\textbackslash\{\}" 2>/dev/null | tail -1 | awk '\{print \textbackslash\{\}\$2\}'); [ -z \textbackslash\{\}"\textbackslash\{\}\$q\textbackslash\{\}" ] \&\& q=none; d=\textbackslash\{\}\$(du -sh \textbackslash\{\}"\textbackslash\{\}\$h\textbackslash\{\}" 2>/dev/null | cut -f1); c=\textbackslash\{\}\$(stat -c \%y \textbackslash\{\}"\textbackslash\{\}\$h\textbackslash\{\}" 2>/dev/null | cut -d' ' -f1 || echo N/A); f=\textbackslash\{\}\$(find \textbackslash\{\}"\textbackslash\{\}\$h\textbackslash\{\}" -type f 2>/dev/null | wc -l); l=\textbackslash\{\}\$(lastlog -u \textbackslash\{\}"\textbackslash\{\}\$u\textbackslash\{\}" 2>/dev/null | tail -1 | awk '\{print \textbackslash\{\}\$5,\textbackslash\{\}\$6,\textbackslash\{\}\$7\}'); [ -z \textbackslash\{\}"\textbackslash\{\}\$l\textbackslash\{\}" ] \&\& l=Never; printf '\%s\textbackslash\{\}t\%s\textbackslash\{\}t\%s\textbackslash\{\}t\%s\textbackslash\{\}t\%s\textbackslash\{\}t\%s\textbackslash\{\}n' \textbackslash\{\}"\textbackslash\{\}\$u\textbackslash\{\}" \textbackslash\{\}"\textbackslash\{\}\$q\textbackslash\{\}" \textbackslash\{\}"\textbackslash\{\}\$d\textbackslash\{\}" \textbackslash\{\}"\textbackslash\{\}\$c\textbackslash\{\}" \textbackslash\{\}"\textbackslash\{\}\$f\textbackslash\{\}" \textbackslash\{\}"\textbackslash\{\}\$l\textbackslash\{\}"; done | sort -t \$'\textbackslash\{\}t' -k3 -hr" | column -t -s \$'\textbackslash\{\}t'} - Sort by disk usage descending, exclude home path and last command
  \item \texttt{sudo bash -c "printf 'USER\textbackslash\{\}tQUOTA\textbackslash\{\}tDISK\_USAGE\textbackslash\{\}tFILES\textbackslash\{\}n'; getent passwd | awk -F: '\textbackslash\{\}\$3 >= 1000 \{ if (\textbackslash\{\}\$1 \textasciitilde{} /\textasciicircum{}nobody\$/) next; print \textbackslash\{\}\$1,\textbackslash\{\}\$6 \}' | while read u h; do q=\textbackslash\{\}\$(quota -u \textbackslash\{\}"\textbackslash\{\}\$u\textbackslash\{\}" 2>/dev/null | tail -1 | awk '\{print \textbackslash\{\}\$2\}'); [ -z \textbackslash\{\}"\textbackslash\{\}\$q\textbackslash\{\}" ] \&\& q=none; d=\textbackslash\{\}\$(du -sh \textbackslash\{\}"\textbackslash\{\}\$h\textbackslash\{\}" 2>/dev/null | cut -f1); f=\textbackslash\{\}\$(find \textbackslash\{\}"\textbackslash\{\}\$h\textbackslash\{\}" -type f 2>/dev/null | wc -l); printf '\%s\textbackslash\{\}t\%s\textbackslash\{\}t\%s\textbackslash\{\}t\%s\textbackslash\{\}n' \textbackslash\{\}"\textbackslash\{\}\$u\textbackslash\{\}" \textbackslash\{\}"\textbackslash\{\}\$q\textbackslash\{\}" \textbackslash\{\}"\textbackslash\{\}\$d\textbackslash\{\}" \textbackslash\{\}"\textbackslash\{\}\$f\textbackslash\{\}"; done" | column -t -s \$'\textbackslash\{\}t'} - Simplified report showing only quota, disk usage, and file count
\end{enumerate}

\section*{Output}
\begin{lstlisting}
USER        QUOTA HOME            DISK_USAGE CREATED    FILES LAST_LOGIN      LAST_CMD
marcos      5G    /home/marcos    290G       2024-01-15 124356   May  20 13:45     grep -r pattern .
libvirt     none  /var/lib/libvirt 16K       2023-06-01 3245    Never              N/A
test-user   2G    /home/test-user 1.2G      2025-03-10 45678    May  10 09:30      ls -lah
admin       10G   /home/admin     3.5G      2024-11-22 89234    May  15 11:22      sudo systemctl status
\end{lstlisting}

\section*{Notes}
\begin{itemize}
  \item This command is designed to run with sudo for complete information; running without sudo will show incomplete results
  \item The 'sudo bash -c' wrapper preserves the entire command context and allows proper variable expansion
  \item Disk quota requires quota system to be enabled on the filesystem; shows 'none' if unavailable
  \item File count calculation uses 'find' which can be slow on directories with millions of files
  \item Account creation date uses home directory birth time (\%w) if available, falling back to modification time (\%y)
  \item Last command is truncated to 30 characters; extracted from .bash\_history. Note: Bash only flushes history to disk upon session exit or when 'history -a' is run, so this shows the last saved command, not necessarily the active session's most recent command.
  \item Users with UIDs < 1000 are system accounts and are excluded by default
  \item The command uses getent instead of /etc/passwd for better compatibility with different user databases (LDAP, NIS, etc.)
\end{itemize}

\section*{Warnings}
\begin{itemize}
  \item Requires sudo access to function properly; user must have sudo privileges without password prompt configured in sudoers for unattended execution
  \item Large file counts or extensive directory trees will cause significant performance degradation
  \item Running with find on large home directories can consume considerable IO and CPU resources
  \item The lslogins command retrieves last login information depending on wtmp/btmp logs; accuracy varies by distribution logging policies
  \item Last command information depends on .bash\_history being available and enabled; other shells use different history files
  \item Network-mounted home directories (NFS, SMB) will significantly increase command execution time
  \item On systems with many users, overall execution time can be substantial; consider running during off-peak hours
\end{itemize}

\section*{See Also}
\begin{itemize}
  \item find-largest-storage-users
  \item du-sort-largest-directories
  \item disk-space-usage-per-directory
  \item find-large-files-recursive
\end{itemize}

\section*{Status}
reviewed

\section*{Safety}
caution

\section*{Shell}
bash

\section*{Platforms}
linux, gnu-linux

\section*{Created At}
2026-04-21

\section*{Updated At}
2026-04-21

\chapter{Git}

\section{Init Repo With Gitignore}
\section*{Initialize git repository, stage all files and create empty .gitignore}
\textbf{Author:} marcos \hfill \textbf{Date:} 2026-04-25

\noindent Initializes a new git repository in current directory, stages all existing files, and creates an empty .gitignore file.

\section*{Language}
git

\section*{Category}
version-control

\section*{Command}
\begin{lstlisting}
git init && git add . && touch .gitignore
\end{lstlisting}

\section*{Explanation}
This command sequence creates a new git repository (git init), adds all files in the current directory to the staging area (git add .), and creates an empty .gitignore file (touch .gitignore) for future ignore patterns.

\section*{Tags}
git, version-control, repository, init, gitignore

\section*{Dependencies}
git


\section*{Arguments}
\begin{enumerate}
  \item \textbf{DIRECTORY} (Optional): Directory where to initialize repository \\ Default: .
\end{enumerate}

\section*{Examples}
\begin{enumerate}
  \item \texttt{git init \&\& git add . \&\& touch .gitignore} - Initialize repository in current directory
  \item \texttt{cd /path/to/project \&\& git init \&\& git add . \&\& touch .gitignore} - Initialize repository in specific directory
\end{enumerate}

\section*{Output}
\begin{lstlisting}
Initialized empty Git repository in /home/user/project/.git/
\end{lstlisting}

\section*{Notes}
\begin{itemize}
  \item The .gitignore file is empty after creation; you should add patterns to ignore unnecessary files
  \item If .gitignore already exists, touch command updates its timestamp
  \item Use git status to verify files are staged
\end{itemize}

\section*{Warnings}
\begin{itemize}
  \item git add . stages all files including potentially sensitive data
  \item Review staged files before committing
\end{itemize}

\section*{See Also}
\begin{itemize}
  \item git-commit-initial
  \item git-ignore-patterns
\end{itemize}

\section*{Status}
reviewed

\section*{Safety}
safe

\section*{Shell}
posix

\section*{Platforms}
linux, gnu-linux, freebsd, openbsd, netbsd

\section*{Created At}
2026-04-25

\section*{Updated At}
2026-04-25

\end{document}
